\documentclass[10pt,a4paper]{amsart}
\usepackage[margin=19mm]{geometry}
\usepackage{amsmath,amssymb,mathtools}
\usepackage[scr=rsfs,cal=euler]{mathalfa}
\usepackage{xfrac}
\usepackage[unicode,hidelinks,bookmarks=false]{hyperref}
\newcommand{\ie}{i.e.\ }
\newcommand{\cf}{cf.\ }
\newcommand{\resp}{respectively\ }
\newcommand{\Subsection}{\subsection}
\providecommand{\nobreakdash}{\nobreak\hbox{-}}
\setlength{\emergencystretch}{2em}
\multlinegap0pt
\usepackage{amssymb}
\usepackage{mathtools}
\newtheorem{thm}{Theorem}
\DeclareMathOperator{\Lip}{Lip}
\DeclareMathOperator{\N}{\mathbb{N}}
\title{A note on attaining diameter two properties in Lipschitz-free spaces}
\author{Jaan Kristjan Kaasik}
\date{Source: arXiv:2604.05569v1 (CC BY 4.0)}
\begin{document}
\maketitle

\noindent\emph{Source and licence.} A note on attaining diameter two properties in Lipschitz-free spaces. Version arXiv:2604.05569v1 (CC BY 4.0), distributed by arXiv under the Creative Commons Attribution 4.0 International licence, http://creativecommons.org/licenses/by/4.0/, as recorded in the arXiv licence metadata for this exact version and shown on the abstract page https://arxiv.org/abs/2604.05569v1. Full licence notice: \texttt{licenses/arxiv-2604.05569-CC-BY-4.0.txt}.\par\smallskip

\noindent\textit{Editorial scope.} Counted unit: the article's thm thm:ASD2P\_main (source0.tex lines 215--217). The article's complete original proof of it is delivered with it (source0.tex lines 218--350, one \texttt{proof} environment of 133 lines). No mathematics is written, repaired or re-derived: the statement and the proof are the article's own lines, in the article's own notation, and no retained line is substituted or re-flowed.  No further source-local context is needed: every cross-reference of every retained line resolves inside the two delivered spans.  Nothing was omitted from any retained span, so the package carries no omission record.  No retained line contains a citation, so the wrapper carries no bibliography and no bibliography entry.  No retained line contains a figure environment, an image-inclusion command or drawing code, so no image is delivered and the package declares no asset dependency.  Editorial formatting only, in addition: an AMS \texttt{amsart} preamble that restates the article's own declarations and macros used by the retained lines, namely the environments \texttt{thm} on separate counters, together with the standard packages those lines need.  The retained environments print in this document as Theorem 1 in order of appearance, whereas the article prints the same items with the numbering of its own full document.  The complete native source of the article as distributed in the arXiv e-print for this exact version is bundled unchanged as \texttt{sources/197964/source0.tex}.
\begin{thm}\label{thm:ASD2P_main}
    Let $M$ be a length space. Then $\mathcal{F}(M)$ has the ASD$2$P. 
\end{thm}

\begin{proof}
Assume $M$ is a length space. 
    Let $n\in \N$, $\lambda_1,\ldots,\lambda_n>0$ with $\sum_i\lambda_i=1$,
    let $f_1,\ldots,f_n\in S_{\Lip_0(M)}$, and let $\alpha>0$. For each $i\in \{1,\ldots,n\}$ consider the slice $S_i=S(f_i,\alpha)$. It suffices to find $\mu,\nu\in \sum_{i=1}^n \lambda_iS_i$ so that $\|\mu-\nu\|=2$.


Fix $i\in \{1,\ldots,n\}$ and let 
\[\varepsilon=\frac{\alpha}{18n}.\]  
Find $u_i,v_i\in M$ with $u_i\neq v_i$ such that
$\langle f_i,m_{u_i,v_i}\rangle>1-\varepsilon$. Denote
\[
r_i=\frac{d(u_i,v_i)}{9n}
\]
and assume without loss of generality that $r_1\leq\ldots\leq r_n$. 
Since $M$ is length, we choose a path $\gamma_i:[0,1]\to M$ joining $u_i$ to $v_i$
such that
\[L(\gamma_i)<(1+\varepsilon)d(u_i,v_i)=\alpha r_i+(1-\varepsilon)d(u_i,v_i),\]
implying
\begin{equation}
\frac{L(\gamma_i)-(1-\varepsilon)d(u_i,v_i)}{r_i}<\alpha.
\end{equation}


Next, we inductively construct points $x_i,y_i\in\mathrm{Im}\,\gamma_i$
so that
\begin{equation}
\min\{d(x_i,x_j),d(x_i,y_j),d(y_i,y_j):\, i>j\}\ge 2r_i
\end{equation}

and

\begin{equation}\tag{2.3}
\min\{d(x_i,v_i),d(y_i,u_i),d(x_i,y_i):\, i=1,\dots,n\}\ge 2r.
\end{equation}

For $i=1$ we set $x_1=u_1$ and $y_1=v_1$.

Let $i\in \{2,\ldots,n\}$ and assume that we have defined points $x_j,y_j\in M$ for every $j<i$. 
To choose $x_i$, we consider the set
\[
F_i=\{x_1,y_1,\dots,x_{i-1},y_{i-1},v_i\}.
\]
We claim that
\[
\mathrm{Im}\,\gamma_i
\setminus\bigcup_{p\in F_i} B(p,2r_i)
\neq\emptyset,
\]
Indeed, if $\mathrm{Im}\,\gamma_i\subset \bigcup_{p\in F_i}B(p,2r_i)$, then
\begin{align*}
 L(\gamma_i)
\le \sum_{p\in F_i}\mathrm{diam}(B(p,2r_i))
\le |F_i|\cdot 4r_i
\le
8nr_i
= \frac89 d(u_i,v_i),
\end{align*}
contradicting $L(\gamma_i)\geq d(u_i,v_i)$.
Thus we may choose $x_i\in\mathrm{Im}\,\gamma_i$
with $d(x_i,F_i)\ge 2r_i$.

Next we define
\[
G_i=\{x_1,y_1,\dots,x_{i-1},y_{i-1},x_i,u_i\}.
\]
The same argument shows that
\[
\mathrm{Im}\,\gamma_i
\setminus\bigcup_{p\in G_i} B(p,2r_i)
\neq\emptyset,
\]
so we may choose $y_i\in\mathrm{Im}\,\gamma_i$
with $d(y_i,G_i)\ge 2r_i$.
This completes the induction and establishes (2.2) and (2.3).

Fix $i\in \{1,\ldots,n\}$. 
By $(2.3)$, we can choose points $z_i,w_i\in\mathrm{Im}\,\gamma_i$
such that
\[
d(x_i,z_i)=r_i,
\qquad
d(z_i,v_i)<d(x_i,v_i),
\]
and
\[
d(y_i,w_i)=r_i,\qquad d(w_i,u_i)<d(y_i,u_i).\]
Applying Lemma~2.1 to the path $\gamma_i$
(with $x=u_i$, $y=v_i$, and $r_i$ as above),
and using (2.1), we obtain
\[
\langle f_i,m_{x_i,z_i}\rangle>1-\alpha
\quad\text{and}\quad
\langle f_i,m_{w_i,y_i}\rangle>1-\alpha.
\]
Hence $m_{x_i,z_i},\, m_{w_i,y_i}\in S_i.$

Now we are ready to define
\[
\mu=\sum_{i=1}^n\lambda_i m_{x_i,z_i}
\quad\text{and}\quad
\nu=\sum_{i=1}^n\lambda_i m_{w_i,y_i}.
\]
Then $\mu,\nu\in\sum_{i=1}^n\lambda_i S_i$.

Define $g:M\to\mathbb R$ by
\[
g(p)=\max\bigl\{ r_i-d(x_i,p),\ r_i-d(y_i,p): i=1,\dots,n\bigr\}.
\]
Each function $p\mapsto r_i-d(x_i,p)$ and $p\mapsto r_i-d(y_i,p)$
is $1$-Lipschitz, hence so is $g$. Set $\tilde g(p)=g(p)-g(0)$. Then 
$\tilde g\in S_{\operatorname{Lip}_0(M)}$.

By (2.2) and (2.3), all of the balls $B(x_i,r)$ and $B(y_i,r)$ are pairwise disjoint.
Moreover,
\[
g(x_i)=r_i,\quad g(z_i)=0,
\quad
g(y_i)=r_i,\quad g(w_i)=0.
\]
Therefore,
\[
\langle \tilde g,m_{x_i,z_i}\rangle=1,
\qquad
\langle \tilde g,m_{w_i,y_i}\rangle=-1,
\]
and consequently,
\[\|\mu-\nu\|\geq
\langle \tilde g,\mu-\nu\rangle
=\sum_{i=1}^n\lambda_i(1-(-1))=2.
\]

 
\end{proof}

\end{document}
